\chapter*{How to read this book}
\addcontentsline{toc}{chapter}{How to read this book}
\markboth{How to read this book}{}

This book has one goal: that you could rebuild IncidentFlow \emph{by hand},
without help, having read it.

That is a higher bar than ``understand the code''. It means you need to know
not only what each file contains, but why it exists, what would break without
it, which command creates it, and in what order the pieces have to be built.

\section*{What is in here}

\begin{description}[leftmargin=2.2cm,style=nextline]
\item[Part I] \textbf{Understanding.} What the system does, the shape of it,
  and where every file lives. Read this once, slowly. Nothing later makes sense
  without it.
\item[Part II] \textbf{The guided tour.} Every screen, as a picture, next to the
  code that produces it and the file path it lives at. This is the bulk of the
  book.
\item[Part III] \textbf{Running it.} Every command, every connection, every
  port. Written for someone who has only ever opened a folder in VS Code.
\item[Part IV] \textbf{Building it from nothing.} The order in which the pieces
  have to be created, and why that order is forced.
\item[Part V] \textbf{Bugs.} Real failures from building this system --- the
  symptom, the wrong guess, the evidence, the fix. This is the part that
  teaches judgement.
\end{description}

\section*{How the pages are laid out}

Three things repeat throughout, and recognising them will make the book faster
to read.

\where{api/app/Models/Incident.php}
\begin{lstlisting}[language=PHP]
// A grey bar like the one above always precedes code, and always
// names the exact file the code came from, relative to the project root.
\end{lstlisting}

\note{A blue box explains \emph{why} something is the way it is. The code shows
what happens; these boxes explain the decision behind it.}

\warnbox{An amber box is a mistake that is easy to make and unpleasant to
diagnose. Read these even if you skim everything else.}

\expect{A green box tells you what you should observe on your own machine if you
followed along correctly. If you see something different, stop --- do not
continue until it matches.}

\bugbox{A red box is a real failure that happened while building this system.
Part V covers them in full, but they appear inline where the relevant code is
introduced, so you meet the bug in the place it lives.}

\section*{The one rule}

Type the commands. Do not copy and paste them, and do not read past a green box
whose result you have not seen on your own screen.

The difference between reading about a system and being able to build one is
almost entirely made up of the small failures you hit while typing --- a wrong
directory, a missing container, a typo in a hostname. Skipping them feels
faster and teaches nothing.
